
\subsection{算子 \texorpdfstring{$\mathfrak q$、$\mathfrak r$}{q、r} 与同伦等价}
\label{subsec:5-3-6}
\liftIndex{I202}\liftIndex{I203}

在第 3.8 节，我们通过插入内部标记点，
推广了此前的构造，并引入
$(C(L;\Lambda_{0,\mathrm{nov}}),
\mathfrak m^{(\mathfrak b,0)})$
这一
$(C(L;\Lambda_{0,\mathrm{nov}}),\mathfrak m)$
的另一种形变；
它由环境辛流形上的闭上链
$\mathfrak b\in
C^2(M;\Lambda^+_{0,\mathrm{nov}})$
参数化。
在第 3.8.7 小节，我们利用这一形变，
在较弱的假设下定义了弗洛尔上同调。
更具体地，当
$L^{(1)}\cap L^{(0)}$
干净相交时，我们定义了
左
$(C(L^{(1)};\Lambda_{0,\mathrm{nov}}),
\mathfrak m^{(\mathfrak b,0)})$、
右
$(C(L^{(0)};\Lambda_{0,\mathrm{nov}}),
\mathfrak m^{(\mathfrak b,0)})$
滤过 $A_\infty$ 双模
\[
\bigl(
C(L^{(1)},L^{(0)};\Lambda_{0,\mathrm{nov}}),
\mathfrak n^{\mathfrak b}_{k_1,k_0}
\bigr).
\]
第 4.6 节证明了：
形变滤过 $A_\infty$ 代数
$(C(L;\Lambda_{0,\mathrm{nov}}),
\mathfrak m^{(\mathfrak b,0)})$
的同伦类不依赖各种辅助选择，
而在 $A_\infty$ 同构的意义下
只依赖于 $L$ 和 $\mathfrak b$ 的上同调类。
本小节要证明，上述定义在
$\Lambda_{\mathrm{nov}}$ 上的
滤过 $A_\infty$ 双模的同伦类
也不依赖这些选择。
这里又有两种情形：
以同一个辛微分同胚移动
两个拉格朗日子流形，
或分别以两条哈密顿同痕移动它们。

设 $(L^{(1)},L^{(0)})$
是 $M$ 中的一对相对 spin
拉格朗日子流形，
两者干净相交。
取辛微分同胚
$\psi\colon M\to M'$，
并记
$L^{(i)\prime}=\psi(L^{(i)})$。
对
$[\mathfrak b]\in
H^2(M;\Lambda^+_{0,\mathrm{nov}})$，
令
\[
\psi^{\mathfrak b}_{*;(i)}\colon
\bigl(
C(L^{(i)};\Lambda_{0,\mathrm{nov}}),
\mathfrak m^{\mathfrak b}
\bigr)
\longrightarrow
\bigl(
C(L^{(i)\prime};\Lambda_{0,\mathrm{nov}}),
\mathfrak m^{\psi_*\mathfrak b}
\bigr)
\]
为定理 4.6.52 的同伦等价。

下面的定理可像定理 5.3.1 一样证明：
\translatornote{原书实体页 335
将下面典范推前的双模结构写作
$\mathfrak m^{\psi_*\mathfrak b}$，
与同段另一实现及定理 5.3.64 的
$\mathfrak n^{\psi_*\mathfrak b}_{k_1,k_0}$
不一致。这里按双模型别统一；
见疑误 A-0310。}
先考虑由
$\mathcal X_1(L^{(i)})$、
$J=\{J_t\}$、
$\mathfrak s_i$ 及 $\mathfrak b$
沿 $\psi\colon M\to M'$
的典范推前所给出的复形
\[
\bigl(
C(L^{(1)\prime},L^{(0)\prime};
\Lambda_{0,\mathrm{nov}}),
\mathfrak n_{k_1,k_0}^{\psi_*\mathfrak b}
\bigr),
\]
再证明这个推前版本与使用另一组选择
$\mathcal X_2(L^{(i)\prime})$、
$J'=\{J'_t\}$、
$\mathfrak s'_i$
所实现的
$(C(L^{(1)\prime},L^{(0)\prime};
\Lambda_{0,\mathrm{nov}}),
\mathfrak n_{k_1,k_0}^{\psi_*\mathfrak b})$
同伦等价。
后一步是下面定理 5.3.65
所处理情形的特例，因此略去证明。

\begin{theorem}{5.3.64}\label{thm:5-3-64}
存在滤过 $A_\infty$ 双模的同伦等价
\[
\begin{aligned}
\varphi_*^{\mathfrak b}\colon\
&\bigl(
C(L^{(1)},L^{(0)};\Lambda_{0,\mathrm{nov}}),
\mathfrak n^{\mathfrak b}_{k_1,k_0}
\bigr)\\
&\longrightarrow
\bigl(
C(L^{(1)\prime},L^{(0)\prime};
\Lambda_{0,\mathrm{nov}}),
\mathfrak n^{\psi_*\mathfrak b}_{k_1,k_0}
\bigr),
\end{aligned}
\]
它覆盖同态对
$(\psi^{\mathfrak b}_{*;(1)},
\psi^{\mathfrak b}_{*;(0)})$。
在下文说明的意义下，
它在同伦意义下不依赖辛同痕的选择，
并在定理 5.3.12 的意义下具有函子性。
\end{theorem}

下面说明定理 5.3.64 中
$\varphi_*^{\mathfrak b}$
在同伦意义下的不变性。
\translatornote{原书实体页 335
在此突然使用未定义的单个
$(M,L)\to(M',L')$ 及
$\psi^t(L)=L'$。
这里按本段的双模两侧
分别写出保持条件；见疑误 A-0309。}
设 $\psi'$ 是另一辛微分同胚
$M\to M'$，
并有一族辛微分同胚
$\{\psi^t\}_{t\in[0,1]}$
把 $\psi$ 与 $\psi'$ 连接起来，
且对 $i=0,1$ 均有
$\psi^t(L^{(i)})=L^{(i)\prime}$。
令
$\varphi_*^{\prime\mathfrak b}$
为与 $\psi'$ 关联的滤过
$A_\infty$ 双模同态。
由于
$\psi_*\mathfrak b$
与 $\psi'_*\mathfrak b$
同调，存在同伦等价
\[
\begin{aligned}
I^{(i)}\colon\
&\bigl(
C(L^{(i)\prime};\Lambda_{0,\mathrm{nov}}),
\mathfrak m^{\psi_*\mathfrak b}
\bigr)\\
&\longrightarrow
\bigl(
C(L^{(i)\prime};\Lambda_{0,\mathrm{nov}}),
\mathfrak m^{\psi'_*\mathfrak b}
\bigr),
\end{aligned}
\]
使
$I^{(i)}\circ\psi^{\mathfrak b}_{*;(i)}$
与
$\psi'{}^{\mathfrak b}_{*;(i)}$
同伦。
此外，存在覆盖
$(I^{(1)},I^{(0)})$
的滤过 $A_\infty$ 双模典范同伦等价
\[
\begin{aligned}
I\colon\
&\bigl(
C(L^{(1)\prime},L^{(0)\prime};
\Lambda_{0,\mathrm{nov}}),
\mathfrak n_{k_1,k_0}^{\psi_*\mathfrak b}
\bigr)\\
&\longrightarrow
\bigl(
C(L^{(1)\prime},L^{(0)\prime};
\Lambda_{0,\mathrm{nov}}),
\mathfrak n_{k_1,k_0}^{\psi'_*\mathfrak b}
\bigr).
\end{aligned}
\]
这里特别指出，
$I$ 与 $I^{(i)}$
完全由
$\psi_*\mathfrak b$
和 $\psi'_*\mathfrak b$ 同调
这一事实经代数方法构造；
它们的同伦类不依赖所涉的各种选择。
于是，
$\varphi_*^{\mathfrak b}$
在同伦意义下的不变性是指：
复合
\[
\begin{aligned}
I\circ\varphi_*^{\mathfrak b}\colon\
&\bigl(
C(L^{(1)},L^{(0)};\Lambda_{0,\mathrm{nov}}),
\mathfrak n^{\mathfrak b}_{k_1,k_0}
\bigr)\\
&\longrightarrow
\bigl(
C(L^{(1)\prime},L^{(0)\prime};
\Lambda_{0,\mathrm{nov}}),
\mathfrak n^{\psi'_*\mathfrak b}_{k_1,k_0}
\bigr)
\end{aligned}
\]
与 $\varphi_*^{\prime\mathfrak b}$
同伦，而它们的覆盖同态对之间
采用上文所述从
$(\psi'{}^{\mathfrak b}_{*;(1)},
\psi'{}^{\mathfrak b}_{*;(0)})$
到
$(I^{(1)}\circ\psi^{\mathfrak b}_{*;(1)},
I^{(0)}\circ\psi^{\mathfrak b}_{*;(0)})$
的同伦。

接着取 $M$ 上两条哈密顿同痕
$\{\psi_\rho^{(i)}\}_\rho$
（$i=0,1$），满足
$\psi_0^{(i)}=\mathrm{id}$，
并置
\[
L^{(0)\prime}
=\psi_1^{(0)}(L^{(0)}),
\qquad
L^{(1)\prime}
=\psi_1^{(1)}(L^{(1)}).
\]
令
\[
\begin{aligned}
{\psi^{(i)}}_*^{\mathfrak b}\colon\
&\bigl(
C(L^{(i)};\Lambda_{0,\mathrm{nov}}),
\mathfrak m^{\mathfrak b}
\bigr)\\
&\longrightarrow
\bigl(
C(L^{(i)\prime};\Lambda_{0,\mathrm{nov}}),
\mathfrak m^{\psi_*^{(i)}\mathfrak b}
\bigr)
\end{aligned}
\]
为定理 4.6.52 构造的同伦等价。
由于
$\psi_*^{(i)}\mathfrak b$
与 $\mathfrak b$ 同调，
定理 4.6.47 保证存在同伦等价
\[
\begin{aligned}
I_{(i)}\colon\
&\bigl(
C(L^{(i)\prime};\Lambda_{0,\mathrm{nov}}),
\mathfrak m^{\psi_*^{(i)}\mathfrak b}
\bigr)\\
&\longrightarrow
\bigl(
C(L^{(i)\prime};\Lambda_{0,\mathrm{nov}}),
\mathfrak m^{\mathfrak b}
\bigr).
\end{aligned}
\]

\begin{theorem}{5.3.65}\label{thm:5-3-65}
存在一个弱滤过 $A_\infty$ 双模同态，
它是下列滤过 $A_\infty$ 双模的同伦等价：
\[
\begin{aligned}
\varphi_*^{\mathfrak b}\colon\
&\bigl(
C(L^{(1)},L^{(0)};\Lambda_{\mathrm{nov}}),
\mathfrak n^{\mathfrak b}_{k_1,k_0}
\bigr)\\
&\longrightarrow
\bigl(
C(L^{(1)\prime},L^{(0)\prime};
\Lambda_{\mathrm{nov}}),
\mathfrak n^{\mathfrak b}_{k_1,k_0}
\bigr).
\end{aligned}
\]
它覆盖同态对
\[
\bigl(
I_{(1)}\circ{\psi^{(1)}}_*^{\mathfrak b},
I_{(0)}\circ{\psi^{(0)}}_*^{\mathfrak b}
\bigr).
\]
它在与上文类似的意义下
不依赖哈密顿同痕的选择，
并在定理 5.3.39 的意义下
具有函子性。
\end{theorem}

定理 5.3.64 与 5.3.65 的证明，
是将本节和第 4.6.6 小节讨论的
论证直接结合。
具体而言，
在构造 $A_\infty$ 双模同态及
同态之间同伦的模空间中加入内部标记点，
再要求这些内部标记点的像落在
$\mathfrak b$ 上，
由此对模空间取截。
细节从略；另见第 7.4.7--7.4.9 小节。
